% ---
% id: protocol-compiler/thm:gates
% kind: thm
% title: "The emission gates and their outcomes"
% summary: >-
%   Each gate is decidable or a presence test, and each failure has one fixed
%   outcome, a refusal or an emission flagged as backing no key length, never a
%   default, dropped entry, changed source mode or licence the compiler writes.
%   Passing asserts nothing about the apparatus or attack-set inclusion.
% prominence: main
% status: implemented
% uses: [thm:validator-soundness, def:eligible, def:gate-outcome-table, def:spec, def:compile, def:homog, def:declaration, def:exact, transmitter/thm:certificate-sufficiency, transmitter/def:source-spec]
% mentions: [prop:conservative, def:validator, rem:gate-by-certificate-form, rem:party-local-soundness, req:gate-outcomes, transmitter/def:srcert, transmitter/fact:sr, transmitter/lem:ensemble-presentation]
% assumes: [A-ATK, A-RHOM]
% refs: []
% provenance: newly-derived
% implements: certiqs_key.protocol_compiler.gates:resolve_gates
% lean:
% notes: "A-RHOM is declared here because gate (G2) decides the certificate once, for the one declaration that describes every round. A-LINK, A-SEQ and A-SRC are not declared: none of them occurs in this theorem's mathematics, and A-SEQ and A-SRC reach its closure through transmitter/thm:certificate-sufficiency. A-ATK's licensed class is carried by its scope (a non-empty list licensed for the declared source mode) and by the statements of its meat-core declarers; this theorem states no such hypothesis, because gate (G3)(a) decides the class."
% ---

One principle fixes every outcome. A check that fails, so that no declaration
is shown to be well-formed, or a certificate that does not verify is a
refusal, and no declaration is emitted; a \tref{well-formed declaration} on
verified data whose key-length claim is unsupported is emitted, not
key-length-eligible, and may be explored numerically. The refusal strings and
the flag that carry these outcomes are Requirement~\ref{req:gate-outcomes}'s.

\begin{theorem}\label{thm:gates}
\begin{enumerate}[label=(\roman*),leftmargin=2.2em]
\item \emph{(decidability)} (G1) is decidable by Theorem~\ref{thm:validator-soundness}; (G2)'s
conditions --- the five of Definition~\ref{def:eligible}\,(G2), namely $\Psi$
Hermitian, $\Psi^2=\Psi$, $\tr\Psi=1$ with the trace rounding nowhere, (SR1)
and (SR2) --- are finitely many exact identities on the emitted data
(\xref{transmitter/thm:certificate-sufficiency}\,(i)); (G3)(a) is decided by
\emph{reading the type of the datum each entry's right-hand-side slot holds,
which is its tag}, and by nothing else --- it compares no matrix against any
other and constructs no right-hand side; (G3)(b) is a presence test. (G2)'s
clause (SR3) and (G3)(b)'s record are attestations and are \emph{not}
decidable; they are tested for presence only;
\item \emph{(outcome of a failure)} If any gate fails, the compiler emits no
key-length-eligible declaration, and which of the \emph{two} permitted
outcomes it produces --- refusal, or emission of a declaration that is not
key-length-eligible --- is fixed by the \tref{gate outcome table}
(Definition~\ref{def:gate-outcome-table}) and is not an
implementer's choice; where the failures of one declaration meet the premises
of a refusal row and of an emission row both, the outcome is the refusal. No
protocol specification meets the premise of rows 2 and 4, the two rows of the
gate outcome table that describe a certificate-absent case: by
Definition~\ref{def:spec}\,(a) a transmitter artefact in trusted-source mode at
$d_A\ge2$ carries a source-replacement certificate. The compiler
never drops a constraint entry silently, never
substitutes a default, never changes the source mode it was handed, and never
synthesizes a licence.
\item \emph{(no assertion)} Passing (G3)(b) records that an attestation was
supplied. It asserts nothing about attack-set inclusion, and no claim of this
contract does. Likewise passing (G2) records that the transmitter's licence
record is present and that its five decided conditions hold of the emitted
data: it asserts nothing about the apparatus, and what it buys for the run is
bought under Assumption~\ref{ass:rhom} --- the certificate is decided once, for
the one declaration that describes every round, which is where this module
spends round homogeneity (\xref{transmitter/fact:sr} drops the per-round
superscript of the source state under exactly that assumption of the
consumer's).
\end{enumerate}

\emph{The general format.} For a declaration in the general format of
Definition~\ref{def:declaration}\,(DX1)--(DX6) the three parts hold with the
following readings, of which the three above are the case of one Alice sector,
one \tref{measured factor}, \tref{POVM form}, a pure-form certificate and no
party-local entry.
\begin{enumerate}[label=(\roman*$'$),leftmargin=2.6em]
\item \emph{(decidability)} (G1) is decidable by Theorem~\ref{thm:validator-soundness} in its
general-format reading. (G2) keeps its five conditions and is decidable in
\emph{one further declared form} (Definition~\ref{def:eligible}): the five
conditions of a pure-form certificate, as above, a certificate declared
in \tref{block form} being refused at this gate; and, for an \tref{ensemble
presentation}, (E1), (E2), (E3) on the declared weighted ensemble together with
the entry-for-entry identity of the declaration's Alice slots with the
transposes of the declared blocks and of their declared-order sum --- finitely
many exact identities, comparisons of labels and integers, and finitely many
summation decisions, the number fixed by $|\mathbb{T}|$, $|\Xset|$ and the
declared dimensions. \emph{Purity is decided by nothing there}: a presentation
declares no $\Psi$, and purity of its blocks, (SR1) and (SR2) are theorems of
the presentation (\xref{transmitter/lem:ensemble-presentation}\,(i), (ii)), the
second giving every non-zero block as a positive multiple of a rank-one
projector. \emph{At one Alice sector it is not thereby waived}, the defined
$\Psi$ being that one block, of unit trace, hence a rank-one projector of which
the five conditions hold. \emph{At two or more it is waived}, and the gate
rests there on two statements --- rank-one blocks by theorem, which is the
hypothesis under which \xref{transmitter/thm:certificate-sufficiency}\,(ii) holds in both directions,
and that theorem's (iii), which reads purity nowhere --- the first of which a
general block form does not supply. (G3)(a) is decided by reading the type of the
datum each entry's right-hand-side slot holds and, for the third form, the tag
field it arrives with --- three forms being in the licensed class, and still
\emph{nothing computed and nothing compared numerically}. (G3)(b) is a presence
test, and (SR3) is tested for presence in every form.
\item \emph{(no row is added)} The outcome of every failure of the general
format is fixed by a \emph{row the gate outcome table already has}. A
general-format check
of Definition~\ref{def:validator} that fails is (G1) and is row 1, the refusal
naming the check. A certificate that is \emph{present} and fails a condition
of (G2) is row 3: a certificate that does not verify is a defect in the
supplied data, not a weaker claim. \emph{Row 3's premise is read of the
declared form}, which reads the cell and adds no row: for a pure-form
certificate the five conditions the cell enumerates are the conditions, and
the cell is literal; for a certificate declared in block form the condition
it fails is the form itself, (G2) admitting no block-form certificate
(Definition~\ref{def:eligible}), and the cell's premise --- a certificate
present and failing a condition of (G2) --- covers it, whether or not it also
fails $\Psi^2=\Psi$; for an ensemble presentation the conditions are (E1) to (E3) and the
identity of the declaration's Alice slots with the transposed declared blocks,
and the cell's premise --- a certificate present and failing a condition of
(G2) --- covers them, the enumeration in it being the pure form's list and not
a restriction of the premise. \emph{An ensemble presentation accompanying a
POVM-form declaration is a failure of that kind and is refused} (row 3),
because nothing ties the declared
$\{M_x\}$ to the presentation (Definition~\ref{def:eligible}). A list
carrying an entry declared general in trusted-source mode is row 5 and is
emitted, not key-length-eligible; \emph{a list carrying a \tref{party-local}
entry is not that row}, that form being licensed in either mode.
\item \emph{(no assertion)} Item (iii) holds as stated, and passing (G3)(b)
records, rather than asserts, one further item: for a list carrying a
party-local entry, what the attack-set licence record attests is a certificate
of the module that produces the chain, on the scoped factor's own channel
(Definition~\ref{def:homog}), and
\emph{no function of the emitted matrices decides it}, exactly as for the
untagged forms. Passing (G2) on an ensemble presentation likewise records that
the transmitter's record is present and that the conditions of \emph{that
form} hold of the emitted data, and asserts nothing about the apparatus; what
it buys for the run is bought under Assumption~\ref{ass:rhom} as in (iii), the
certificate being decided once for the one declaration that describes every
round.
\end{enumerate}
\end{theorem}

\begin{proof}
\emph{Idea.} Each gate condition is either a finite list of exact identities,
a type read off each constraint entry, or a presence test; the outcome of a
failure is the one the gate outcome table of
Definition~\ref{def:gate-outcome-table} fixes, no row of it changes the
source mode or drops an entry, and by the definition of a protocol
specification no specification meets its certificate-absent rows; and
the general format adds decisions of the same kinds and failures that land on
existing rows.

\emph{Step 1 (decidability, (i)).} Each clause is as stated:
Theorem~\ref{thm:validator-soundness}\,(iv) for (G1);
\xref{transmitter/thm:certificate-sufficiency}\,(i) for the (G2) identities.
For (G3)(a): by Definition~\ref{def:homog} the licensed class is defined by
which of the two forms each entry \emph{declares}, and the tag is a field of
the datum (Definition~\ref{def:declaration}), so deciding the class is reading
$m$ tags and checking that the datum each names has the type it names --- $m$
comparisons of a label and a type, terminating and exact. \emph{Nothing is
computed and nothing is compared numerically.} Presence tests are decidable
trivially. That (SR3) and the licence record are attestations is
\xref{transmitter/def:srcert}\,(SR3) and Assumption~\ref{ass:atk}: neither is
a property of the emitted matrices, so no function of them decides it.

\emph{Step 2 (the marginal slot and the certificate-absent rows, for (ii)).}
$\tau_A$ reaches the emitted
data only through the transmitter artefact's marginal slot ---
Definition~\ref{def:compile} passes that slot through, and no input of this
module carries a marginal or a mode separately (Definition~\ref{def:spec}\,(a):
the mode is read off the artefact's marginal slot and is never a second
datum). The premise of
rows 2 and 4 is an artefact in trusted mode at $d_A\ge2$ whose certificate is
absent, and by Definition~\ref{def:spec}\,(a) a transmitter artefact in
trusted-source mode at $d_A\ge2$ additionally carries a source-replacement
certificate; so no protocol specification meets that premise. That the engine
refuses such a request is Requirement~\ref{req:gate-outcomes}'s, and no step
here uses it. At
$d_A=1$, where (G2) does not demand the certificate at all, the slot holds
$\tau_A=[1]$, the value the transmitter filled it with
(Definition~\ref{def:spec}\,(a)), and this
module passes the slot through --- that is not a default substituted for a
missing datum here, which item (ii) forbids, because this module fills
nothing: $\Seq{A}=\{[1]\}$ is a singleton, so D(ii) admits that value and no
other, and at $d_A=1$ the two source modes impose the same condition
(\xref{transmitter/def:source-spec}), so the filled slot excludes nothing.

\emph{Step 3 ((ii) and (iii)).} (ii) and (iii) are rules of the
\tref{compilation map}, stated here because they are the behaviour the gate exists to produce.
(ii)'s precedence clause is part of that rule and is what makes the table fix
\emph{one} outcome for each declaration: one declaration can meet the premise
of a refusal row and of an emission row at once --- a certificate failing a
condition of (G2), row~3, beside a list carrying an entry declared general,
row~5 --- and the clause takes the refusal, so no declaration whose supplied
data failed a check or a certificate is emitted.
(ii) is consistent: the one non-refusal outcome is exactly the one for which
this contract proves the result sound --- an ineligible declaration is not
consumed by a key-length statement at all --- and no row changes the source
mode, re-reads a constraint entry or drops one. No row invokes
Proposition~\ref{prop:conservative}: the rows on which it could bear, the
certificate-absent ones, are met by no protocol specification (Step~2).
(iii)'s second sentence is
Assumption~\ref{ass:rhom} read where it is spent: (G2) decides the certificate
once and \xref{transmitter/thm:certificate-sufficiency}\,(iii) is a statement
about one round's source state, so the flag says of the run only what one
declaration can say of every round.

\emph{Step 4 (the general format, (i$'$)).} Theorem~\ref{thm:validator-soundness}\,(iv) in its
general-format reading gives (G1). For (G2): each condition
Definition~\ref{def:eligible} lists for a form is either an identity or an
order relation between declared numbers, decided as
Definition~\ref{def:exact} decides such a predicate, or a comparison of labels,
integers and structures, or one summation decision --- and each list is finite,
its length fixed by $|\mathbb{T}|$, $|\Xset|$ and the declared dimensions. That
purity is decided by nothing \emph{for a presentation} is read off
\xref{transmitter/def:srcert}: a presentation declares no $\Psi$, and purity,
(SR1) and (SR2) are conclusions of the transmitter's lemma, so no function of
the emitted data is asked for them. Three separate statements describe what
the gate then demands. \emph{First}, the five conditions are decided of every
certificate declared in pure form, and at one Alice sector --- the only case a
declaration in the \tref{basic format} can present --- they \emph{hold} of the
$\Psi$ a presentation defines, by
\xref{transmitter/lem:ensemble-presentation}\,(ii) with unit \tref{total weight}.
\emph{Second}, a presentation is admitted only beside an
\emph{ensemble-form} declaration, which no basic-format declaration is; and a
certificate declared in \emph{block form} is refused here
(Definition~\ref{def:eligible}). \emph{Third}, at $|\mathbb{T}|\ge2$ the
defined $\Psi=\bigoplus_t\Psi_t$ has $\Psi^2=\bigoplus_tg_t\Psi_t\ne\Psi$ once
two blocks are non-zero, so global purity is not demanded of a presentation
there, and what the gate rests on in its place is what
Remark~\ref{rem:gate-by-certificate-form} names --- rank-one blocks by theorem,
giving both directions of \xref{transmitter/thm:certificate-sufficiency}\,(ii),
and that theorem's (iii), which reads purity nowhere. A general block form
lacks the first and has the second, that theorem being stated in every form;
the refusal of a block form rests on scope and not on that difference, and
the distinction is not the word ``purity''
(Remark~\ref{rem:gate-by-certificate-form}). For (G3)(a): by
Definition~\ref{def:homog} the licensed class is defined by which of the three
forms an entry \emph{declares}, the first two by the type of the datum and the
third by the tag field beside a real scalar, so deciding the class is reading
$m$ tags and checking that the datum each names has the type it names --- $m$
comparisons of a label and a type, terminating and exact, and again nothing
computed: no form of the general format introduces a numerical comparison.

\emph{Step 5 (the general format, (ii$'$)).} (ii$'$) is a statement about
where each failure lands, and it is proved by exhibiting the row: a failing
check is (G1) by the definition of that gate, and (G1)'s row is a refusal
naming the check, so every failing check of the general format lands on that
row; a certificate present and failing a condition of (G2) is the premise of
row 3, its outcome a refusal naming the condition. \emph{A block-form
certificate meets that premise by its form}: the record is present, and the
condition of (G2) it fails is that (G2) admits no certificate declared in
block form (Definition~\ref{def:eligible}), which holds of every block-form
certificate whatever its blocks --- including one whose one non-zero block is
a rank-one projector and which therefore meets $\Psi^2=\Psi$ --- so the
outcome is a refusal naming the form. \emph{An ensemble presentation
beside a POVM-form declaration meets the premise too}: the record is present,
and the condition of its form that it fails is the one requiring the
declaration's Alice slots to be \emph{the} slots the presentation defines,
which a POVM-form declaration has not got --- so the outcome is a refusal, and
the refusal names which condition failed. No premise of any other row is met
by a general-format failure, and each of the six rows keeps its outcome; in
particular no general-format failure yields a key-length-eligible
declaration, and none yields an emission on a row whose outcome is a refusal.

\emph{Step 6 (the general format, (iii$'$)).} The added item is the
\tref{soundness condition of a party-local entry} (Definition~\ref{def:homog};
that the party-local form is sound only under it is
Remark~\ref{rem:party-local-soundness}), read where the record is tested:
every attack reaches the scoped factor through a channel on that factor alone
carrying the spectral certificate, which is a certificate of the module that
produces the chain (Definition~\ref{def:eligible}), and the gate tests for
the presence of the record that attests it. \emph{So admitting the third form
at (G3)(a) moves no decision from an attestation to this module and none from
this module to an attestation}: it adds one licensed tag and one clause to
what the same record attests. The rest of (iii) holds as stated, including
where Assumption~\ref{ass:rhom} is spent.
\end{proof}
